% Revised 2026-11-06: external-tool scope (OBLITERATUS: Qwen2.5-7B only; jwest33: + Qwen3.5-4B); seed caveat with
% the companion study's replicates.
% Revised 2026-11-06 (draft v3): the selection-cap re-export is done (the "needs a GPU run" TODO is resolved); its remaining
% limits stated (one study, one seed, one model, one cap, general direction only; code study not re-exportable because
% its prompt files are missing: reexport_1x/VERSION.md "Inputs not in /data"); external tools no longer pending;
% jwest33 Qwen3.5-4B: as-released run primary, hybrid-detection run sensitivity (author decision); 7B 1.x cells judged.
\subsection{Limitations and Future Work}
\label{sec:limitations}

The evidence is narrower than the question. All base models of the main design are from one family (Qwen2.5), the
like-for-like contrast rests on two base models, and Heretic~2.0 was not run on Qwen2.5-14B or the coder models. The
selection-cap check rests on a single re-exported trial: one search with one seed on one model (Qwen2.5-32B), for the
general direction only and at one cap (dev KL $\le$ 0.1). The 1.x code study could not be re-exported because its
prompt files were not archived with the run records, the Qwen2.5-7B searches and intermediate caps were not re-exported, and the
re-exported arm still differs from 2.0 DiM in trial budget, seed, system prompt, marker list and KL objective. The
external tools were each run once at their documented default or recommended settings, without a search or a
selection cap, so each cell describes one preset of one tool at an uncontrolled divergence rather than the best that
tool can do. OBLITERATUS was evaluated on Qwen2.5-7B only, because it does not support Qwen3.5-4B, and jwest33's
abliterator on Qwen2.5-7B and on Qwen3.5-4B, where the run as released is the primary cell and the run with the
restored hybrid-architecture detection a sensitivity analysis; the Qwen3.5-4B comparison lies outside the main
design and uses the companion study's Heretic~2.0 cells. Both tools' prompt pools partly overlap the held-out set
(\S\ref{sec:external}). ErisForge, llm-abliteration and FailSpy are represented only by the keyword-based results of
the earlier comparison. Each configuration contributed one search or run with one seed per model, so every test
compares two released checkpoints, and differences of the size of the companion study's seed-to-seed variation (a
standard deviation of about 13~pp for the difference between two independently optimised cells
\cite{young2026harder}) cannot be attributed to the tool or version with confidence. Responses were greedy and
limited to 512 tokens, so refusal later in long answers and sampling effects are not captured. Residual refusal
measures willingness, not the usefulness of a harmful answer; the HarmBench classifier's compliance label is only a
proxy \cite{carleo2026willing}. No human audit of the judged responses was done, so every rate is
relative to the automatic judges; an audit of a stratified sample is left to future work. The 1.x code variant's prompt files could not
be recovered (only the size of its dev set, 100 prompts, is recorded), so its overlap with the held-out code prompts was bounded by excluding all 41
prompts from the same corpora rather than checked item by item. The procedures behind the community releases are
undocumented.

Five next steps follow directly. Re-exporting the archived 1.x trials at intermediate caps (dev KL 0.02 and 0.05) on
Qwen2.5-32B and at the 2.0 cap on Qwen2.5-7B, and recovering the code study's prompt files so that its trials can be
re-exported, would show how the held-out effect of the cap scales and whether it holds beyond one search. Running the
2.0 DiM search with the 1.x trial budget, seed, system prompt and marker list, or the reverse, would separate these
settings from the version once the cap is matched. A second search seed per arm, starting with the companion study's
existing seed-2 Qwen2.5-7B cells, would show how much of each contrast exceeds search-to-search variation. Running
the external tools across their documented settings, and at a divergence cap where the tool allows one, would
separate preset from algorithm, and adding at least one non-Qwen family would test generality. Releases of
abliterated models should state the selection rule and the divergence at which they were selected.
